% GPT-6 Astra Pro; effort not exposed; /root/replacement; 22 November 2026.
\section[HD-F1-R30: parallel surface defects cannot cancel]{HD-F1-R30: the actual common-quotient cancellation excludes rank-one divisor factors}
\label{sec:followup-rm-parallel-surface-cancellation}

Every actual common-quotient object in the class
of R21--R29 has rank-two divisor factors in
(F1R28.1). Together with R29, this proves that
both factors are definite in the common-quotient
and three-row common-image branches.
The same actual finite point action and
connecting class are retained throughout.
The proof compares R22's necessary
degree-two image with the Serre transpose of
the actual hull quotient. It allows arbitrary
infinitesimal thickenings and lower-dimensional
parts of that quotient.

\paragraph{1. The actual quotient and its pure surface quotient.}
Suppose the two divisor factors in the
common-quotient branch had rank one.
Use R22's actual coherent exact sequence
\[
       0\longrightarrow T
          \xrightarrow{\iota}L^{\oplus2}
          \xrightarrow{q_Z}Z\longrightarrow0.
 \tag{F1R30.1}
\]
The image of the actual invariant map
\[
 q_Z^*:\operatorname{Ext}^2_G(Z,L)
                \longrightarrow
                    \operatorname{Ext}^2_G(L^{\oplus2},L)
 \tag{F1R30.2}
\]
has dimension at least three by (F1R22.16).
The degree-two cycle of $Z$ is nonzero.
By R29, all its surface components are
translates of the same abelian surface
$B\subset X$.

Let $Z_0\subset Z$ be its largest coherent
subsheaf supported in dimension at most one,
and put $P=Z/Z_0$. This subsheaf exists
because sums preserve the dimension bound
and $Z$ is Noetherian. Its maximality
makes $P$ pure of dimension two: a
nonzero subsheaf of $P$ of lower dimension
would have a preimage in $Z$ of dimension
at most one, contradicting maximality.
Thus the entire set-theoretic support of
$P$ is a finite union of the same
translates of $B$. Its generic multiplicities
are unchanged, and its full scheme-theoretic
structure is kept in every subsequent step.
The removed $Z_0$ is accounted for by the
actual Ext isomorphism below.
The construction is preserved by the
given translations and line tensors,
so $Z_0$ and $P$ inherit the actual lifts.

Because $Z_0$ has codimension at least
three on the smooth fourfold,
\[
       \operatorname{Ext}^1_X(Z_0,L)
          =\operatorname{Ext}^2_X(Z_0,L)=0.
 \tag{F1R30.3}
\]
Indeed all its sheaf Ext groups in degrees
below three vanish, and the local-to-global
Ext spectral sequence gives both displayed
global vanishings. This is the precise
vanishing of Huybrechts--Lehn
Proposition~1.1.6(i), after the actual
line twist. Apply Hom into $L$ to
$0\to Z_0\to Z\xrightarrow{p}P\to0$.
The two neighboring zero groups in
(F1R30.3) give the actual isomorphism
\[
 p^*:\operatorname{Ext}^2_X(P,L)
                    \xrightarrow{\ \sim\ }
                    \operatorname{Ext}^2_X(Z,L).
 \tag{F1R30.4}
\]
It also holds on invariants. If
$q_P=pq_Z:L^{\oplus2}\to P$, composition
gives
\[
                       q_P^*=q_Z^*p^*.
 \tag{F1R30.5}
\]
Hence it suffices to bound the ordinary
rank of $q_P^*$ by two. The invariant
image in (F1R30.2) is contained in that
ordinary image.

\paragraph{2. The actual quotient torus retains every thickening.}
Let
\[
                    \pi:X\longrightarrow Y=X/B
 \tag{F1R30.6}
\]
be the quotient homomorphism; $Y$ is
an abelian surface. On universal covers
it is the projection $V\to V/T_0B$
and on integral homology it is the
associated quotient of lattices.
For completeness its projectivity can
be seen from the given polarization:
the rational polarization-orthogonal
complement of $H_1(B,\mathbf Q)$ gives
a complementary complex subtorus $B'$,
with an actual addition isogeny
$B\times B'\to X$. The induced map
$B'\to Y$ is a quotient by its finite
kernel. Identify its universal cover with
that of $B'$; its lattice contains the
lattice of $B'$ with finite index.
Multiplying the restricted positive
Riemann form by a sufficiently divisible
integer makes it integral on this larger
lattice. The actual Appell--Humbert
construction then gives a positive line
bundle on the quotient. This supplies the
abelian quotient in (F1R30.6);
the differential and lattice maps
are the stated ones.

Put $\mathcal F=P\otimes L^{-1}$
and write
\[
          s=q_P\otimes1_{L^{-1}}:
                      \mathcal O_X^{\oplus2}\longrightarrow\mathcal F.
 \tag{F1R30.7}
\]
The set $\pi(\operatorname{Supp}\mathcal F)$
is a finite set $S\subset Y$.
If $\mathcal I_S$ is its reduced ideal,
coherence and the Noetherian property
give an integer $N$ for which
\[
                   (\pi^{-1}\mathcal I_S\cdot\mathcal O_X)^N
                                      \mathcal F=0.
 \tag{F1R30.8}
\]
Locally this follows from
$\pi^{-1}\mathcal I_S\subset
\sqrt{\operatorname{Ann}\mathcal F}$;
finitely many generators and a finite
affine cover give one common exponent.
Thus the actual $\mathcal F$, with
all its thickening, is supported
scheme-theoretically on
$\pi^{-1}(S_N)$, where $S_N$ is the
finite subscheme defined by
$\mathcal I_S^N$. It has not been
replaced by its restriction to $S_{\rm red}$.

Choose a projective embedding of $Y$
and a hyperplane missing every point
of $S$. Its complement in $Y$ is an
affine open $U$ containing the entire
finite subscheme $S_N$. Let
\[
                         V=\pi^{-1}(U)\subset X.
 \tag{F1R30.9}
\]
The whole sheaf $\mathcal F$ is supported
inside $V$. For every $n$, the actual
restriction map is an isomorphism
\[
 \operatorname{Ext}^n_X(\mathcal F,\mathcal F)
          \xrightarrow{\ \sim\ }
 \operatorname{Ext}^n_V(\mathcal F|_V,\mathcal F|_V).
 \tag{F1R30.10}
\]
Here is a precise support argument.
Every coherent sheaf
$\mathcal Ext^q_X(\mathcal F,\mathcal F)$
is supported on the closed subset
$\operatorname{Supp}\mathcal F\subset V$.
Its cohomology computed on $X$ or
on $V$ is the same: as an abelian
sheaf it is the exact pushforward
from that closed subset, for which
the two global-section functors agree.
Equivalently one can use its actual
annihilator thickening, which is closed
in both spaces. The natural comparison
of local-to-global Ext spectral sequences
is therefore an isomorphism on every
$E_2^{p,q}$ term and on its abutment.
The boundedness needed here follows
from smoothness of $X$ and coherence
of $\mathcal F$. This proves (F1R30.10)
without discarding any normal nilpotent
direction.

\paragraph{3. The base one-forms act by zero on this sheaf.}
Let
\[
 U_0=\pi^*H^1(Y,\mathcal O_Y)
                   \subset H^1(X,\mathcal O_X)=\mathcal H.
 \tag{F1R30.11}
\]
The cohomology space $\mathcal H$ has
dimension four.
The inclusion is injective and
$\dim U_0=2$: under the invariant-form
isomorphism $H^1(\mathcal O)=
\overline{T_0(\,\cdot\,)}^*$ it is the
dual of the surjective conjugate
differential of (F1R30.6).

For $\alpha\in U_0$, view its usual
cohomology class as the actual morphism
$\mathcal O_X\to\mathcal O_X[1]$.
Tensoring with $\mathcal F$ gives
\[
                   \alpha_{\mathcal F}:
                                  \mathcal F\longrightarrow\mathcal F[1].
 \tag{F1R30.12}
\]
Write $\alpha=\pi^*\alpha_Y$.
Since $U$ is affine,
$H^1(U,\mathcal O_U)=0$.
The restriction of $\alpha_Y$ is zero
as a morphism
$\mathcal O_U\to\mathcal O_U[1]$.
Its pullback to $V$ is also zero,
so (F1R30.12) restricts to zero.
The injectivity in (F1R30.10) proves
\[
                         \alpha_{\mathcal F}=0
                                   \quad(\alpha\in U_0).
 \tag{F1R30.13}
\]
This is a nullity statement for the
original thickened sheaf, not only
for its set-theoretic fibers.

Write $s_1,s_2:\mathcal O_X\to\mathcal F$
for the actual columns of (F1R30.7).
For $\alpha\in U_0$ and
$\xi\in\mathcal H$, functoriality of the
tensor action gives the exact composition
\[
 s_i[2](\alpha[1]\xi)
       =\alpha_{\mathcal F}[1]\,s_i[1]\xi=0.
 \tag{F1R30.14}
\]
Thus each of the actual cohomology
maps
\[
              H^2(X,\mathcal O_X)
                         \xrightarrow{(s_i)_*}H^2(X,\mathcal F)
 \tag{F1R30.15}
\]
kills $U_0\wedge\mathcal H$.
The invariant-form exterior description
of an abelian variety gives
$H^2(X,\mathcal O_X)=\bigwedge^2\mathcal H$.
Consequently (F1R30.15) factors through
the actual quotient
\[
 \bigwedge\nolimits^2\mathcal H/(U_0\wedge\mathcal H)
             =\bigwedge\nolimits^2(\mathcal H/U_0),
       \qquad
 \dim\bigwedge\nolimits^2(\mathcal H/U_0)=1.
 \tag{F1R30.16}
\]
There are two actual columns, so
\[
       \operatorname{rank}
       \left(H^2(X,\mathcal O_X)^{\oplus2}
                         \xrightarrow{s_*}H^2(X,\mathcal F)\right)
                                     \le2.
 \tag{F1R30.17}
\]

\paragraph{4. Serre duality compares the same quotient map.}
Fix a nonzero invariant volume form
to identify $\omega_X=\mathcal O_X$.
Serre duality on the smooth fourfold
identifies the transpose of
\[
 q_P^*:\operatorname{Ext}^2_X(P,L)
                  \longrightarrow\operatorname{Ext}^2_X(L^{\oplus2},L)
 \tag{F1R30.18}
\]
with the actual postcomposition map
\[
 (q_P)_*:\operatorname{Ext}^2_X(L,L^{\oplus2})
                   \longrightarrow\operatorname{Ext}^2_X(L,P).
 \tag{F1R30.19}
\]
Explicitly, for $a:P\to L[2]$ and
$b:L\to L^{\oplus2}[2]$, the equality
of their Serre pairings is the
literal equality of compositions
\[
 \operatorname{Tr}_L\!\left(a[2]q_P[2]b\right)
       =\operatorname{Tr}_L\!\left(a[2](q_P[2]b)\right).
 \tag{F1R30.20}
\]
All sheaves here are perfect because
$X$ is smooth; no local-freeness of
$P$ or $\mathcal F$ is required.
Tensoring both arguments by $L^{-1}$
identifies (F1R30.19) with the map
$s_*$ in (F1R30.17).
Thus (F1R30.18) has rank at most two.

The isomorphism (F1R30.4) and
composition (F1R30.5) now bound the
ordinary image of $q_Z^*$ by two,
and hence bound its invariant image
by two. This contradicts the actual
lower bound of three in (F1R30.2).
The common-quotient rank-one-factor
case is excluded.

\paragraph{Exact scope and sources.}
R22 supplies the original map with
its required image of dimension at
least three, and R29 supplies the
actual parallel surface supports.
Huybrechts--Lehn, second edition,
Definition~1.1.4 and Proposition~1.1.6(i),
give the torsion filtration and the
precise low-degree Ext vanishings;
their used construction and complete
proposition proof were freshly read.
The power-annihilator argument is
Stacks Lemma~30.10.2, Tag~01Y9, with
its local proof in Lemma~10.62.4,
Tag~00L6. The exact closed-support
cohomology comparison is Lemma~20.20.1,
Tag~02UV, and the underlying exact
closed pushforward is Lemma~17.6.1,
Tag~01AX. Their stated proofs were
read; the actual Ext restriction and
cup-action maps are given above.

The deletion of $Z_0$ is justified
by both neighboring global Ext
vanishings, and only the resulting
isomorphic source is used in the
same predecessor map. The affine
restriction preserves every finite
thickening over $S$, and the final
rank comparison is ordinary before
taking invariants. No invariant
affine neighborhood, reduced-support
replacement, or vanishing of all
$H^2(X,\mathcal F)$ is assumed.
The remaining definite-factor
common-quotient objects, the
definite-factor three-row objects,
and the two-row common-image class
still require their actual coupling
and full trace.

